\section{Discussion}
\label{sec:discussion}

The contribution of Notebook-to-OSL is best assessed through the three research questions rather than through the workflow components in isolation. The situated session, interview, annotation schema, operator review, and OSL transformation are instruments used to examine distinct uncertainties in the path from operator practice to formal representation.

\subsection{RQ1: Situated training compared with post-session interview}

RQ1 asks how situated CNC operator training differs from a conventional post-session interview in the operator-strategy knowledge it elicits. The proposed design does not assume that situated training is categorically superior. Situated training may reveal perceptual cues, artefact use, timing, contextual dependencies, and corrections that become visible only while work is demonstrated. A post-session interview may instead produce more reflective explanations, broader summaries, or knowledge that the operator did not articulate during action.

The relevant result is therefore not the number of statements produced by each source. The analysis should identify which knowledge elements are shared, situated-only, or interview-only, and whether those differences affect later interpretation or modelling. A meaningful contribution would be evidence that the elicitation setting changes the content or structure of the operator-strategy knowledge available to Digital Twin development. If the two sources produce materially equivalent fragments, the paper must report that result rather than assuming a situated advantage.

\subsection{RQ2: Changes introduced by operator validation}

RQ2 asks how operator validation changes the researcher's interpretation of captured notes. This question addresses a central methodological risk: field material is not transferred directly into an engineering artefact but interpreted by the researcher. The initial annotation may misidentify an observation, omit a condition, collapse alternatives into one action, infer a rationale, or assign responsibility incorrectly.

The validation record should therefore preserve substantive changes rather than only a final approval status. Corrections show where the initial interpretation was wrong; narrowing shows where it was too general; rejection shows where the source did not support the proposed fragment; and sensitivity decisions show where valid knowledge cannot be used as ordinary research material. The frequency and character of these changes provide evidence about the reliability of the transformation process and about which annotation categories require the strongest review.

This question also prevents operator validation from becoming ceremonial member checking. A method that records only accepted fragments would conceal the interpretive work needed to produce them. The research value lies partly in making disagreement and revision visible.

\subsection{RQ3: Transformation into semantically coherent OSL representations}

RQ3 asks to what extent validated fragments can be transformed into traceable, semantically coherent OSL representations and what gaps remain. This question separates formal representation from storage or syntax generation. A fragment is not model-ready merely because its fields can be serialised as JSON or rendered as a code listing. The representation must preserve the distinctions relevant to the validated knowledge, including applicability, observations, interpretations or hypotheses, alternatives, response roles, expected consequences, evidence, uncertainty, provenance, and unresolved criteria.

The transformation analysis should distinguish four outcomes: coherent representation, coherent but explicitly incomplete representation, partial representation with information loss, and unsupported mapping. Failed or incomplete transformations are analytically useful. They may reveal that the empirical fragment lacks a decision criterion, that the annotation schema omitted a relevant distinction, or that the current OSL profile cannot express an important relationship. RQ3 therefore evaluates both the readiness of the fragments and the adequacy of the target representation without claiming operational correctness, runtime recommendation, or safety.

\subsection{Integrated contribution to Digital Twin development}

Together, the questions define a bounded engineering-informatics contribution. RQ1 examines how operator knowledge enters the process. RQ2 examines how that knowledge changes through human review. RQ3 examines whether the reviewed knowledge can support formal engineering representations. The method therefore contributes an auditable chain from elicitation source to interpretation version, validation decision, and representation outcome.

This chain is more important than the note-taking application alone. A simple capture tool has little research value by itself. The value lies in exposing where knowledge was obtained, where the researcher changed it, where the operator corrected it, and where formalisation succeeded or failed. Such traceability can support later requirements work, OSL modelling, explanation design, validation planning, or semantic Digital Twin development.

\subsection{Evaluation limits}

A small pilot can demonstrate feasibility and reveal consequential differences, interpretation changes, and representation gaps. It cannot establish general superiority of situated training, generalise across CNC operations or manufacturing domains, validate OSL as a complete language, or show that a resulting Digital Twin improves production performance. Stronger claims would require multiple operators, tasks, sites, independent analysis, or downstream use studies.

The main threat to RQ1 is contamination between the situated session and the post-session interview. This should be reduced by preserving source separation, using the same task scope, and avoiding situated-note prompts during the initial interview pass. The main threat to RQ2 is social agreement during operator review; disposition reasons and concrete revisions should therefore be recorded. The main threat to RQ3 is circularity, because the annotation schema and OSL may have been developed together. The analysis should preserve unsupported and information-losing mappings rather than adapting every fragment until it appears representable.

Other limitations remain. Note-taking may miss embodied actions, gestures, timing, or sensory detail. Operators may find some knowledge difficult to verbalise. Researcher participation may influence the explanation. Industrial confidentiality may restrict what can be stored, quoted, or published. These limitations should be documented in the reflection and validation records rather than hidden by the final formal representation.

\subsection{Ethical and confidentiality considerations}

Raw notes and interview records may contain proprietary details, customer information, machine programs, identifiable operator statements, or commercially sensitive parameters. The study should distinguish protected empirical material from publishable examples and preserve confidentiality decisions throughout the trace. A fragment marked \texttt{sensitive} may remain methodologically relevant to RQ2 or RQ3 while being excluded from public datasets, listings, and quotations.
